\section{Explicit descent on the coherent domain}\label{sec:finite}

We now fix a parameter tuple and a finite domain independently of the preceding local continuation hypotheses:
\begin{equation}
\eps=\frac1{10},\qquad g=\frac1{60},\qquad
\lambda=\frac1{300},\qquad
k=\left(\frac{999}{1000},1,\frac{1001}{1000}\right).
\label{eq:parameters}
\end{equation}
With \(D_0=\pi/12\), let
\begin{equation}
\D=\left\{\Omega\in\C^3:
\frac9{10}\le|\Omega_i|\le\frac{11}{10},\
\text{ordinary phase width}\le D_0\right\}.
\label{eq:domain}
\end{equation}
The containing phase interval can have any common orientation. It is an ordinary full-turn interval, not phase modulo \(2\pi/3\).

We use the following precise inherited facts from \cite{relaxation}, Sections 6--7, \texttt{lem:phase}, \texttt{eq:lower}, \texttt{eq:upper}, and \texttt{thm:convergence}:
\begin{enumerate}
\item Positive prestage mixing keeps its phases in the input interval, and
\begin{equation}
\frac{88701}{100000}+\frac3{100}\cos D_0
\le |v_i|\le\frac{107701}{100000},\qquad
\frac9{10}<|v_i|<\frac{11}{10}.
\label{eq:prebounds}
\end{equation}
\item For positive-modulus prestages of ordinary phase width \(D_v<\pi/6\) and \(0\le\lambda\le1/9\), the simultaneous phase stage preserves its interval and satisfies
\begin{equation}
|\psi_i^+-\psi_j^+|
\le(1-9\lambda\cos3D_v)|\psi_i-\psi_j|.
\label{eq:phasebound}
\end{equation}
\item The domain \(\D\) is forward invariant. There is a unique aligned fixed amplitude profile \(r_*\in[9/10,11/10]^3\), satisfying \eqref{eq:fixedbranch} at \eqref{eq:parameters}. Its common rotations are fixed states.
\end{enumerate}
These inputs include positive \emph{prestage} bounds; positive input amplitudes alone would not justify the regularity needed later. No profile value is numerically computed.

\begin{theorem}[Finite-domain native descent]\label{thm:finite}
For every \(\Omega\in\D\) at \eqref{eq:parameters},
\begin{equation}
\boxed{\V(F\Omega)-\V(\Omega)
\le-\frac{16869}{20000}\|v-\Omega\|_2^2.}
\label{eq:finitedescent}
\end{equation}
Equality in monotonicity, \(\V(F\Omega)=\V(\Omega)\), holds exactly on
\(\{e^{i\theta}r_*:\theta\in\R\}\). The coefficient is a valid bound, without an optimality claim, and multiplies the squared prestage increment.
\end{theorem}

\subsection{A whole-segment curvature estimate}
\begin{proof}
Both \(\Omega\) and \(v\) lie in the convex product of closed radius-\(11/10\) disks. For \(b=v-\Omega\),
\begin{equation}
|\Omega_i+tb_i|
\le(1-t)|\Omega_i|+t|v_i|\le11/10,\qquad 0\le t\le1.
\label{eq:segment}
\end{equation}
No convexity of the annulus is assumed. The potential is polynomial in real coordinates, so this segment argument needs no nonzero-phase condition at its interior points.

The onsite Hessian at \(x_i\in\R^2\) is
\begin{equation}
\eps\big[(|x_i|^2-k_i)I_2+2x_ix_i^T\big],
\label{eq:onsitehessian}
\end{equation}
and the graph Hessian is \(-gL\otimes I_2\). For any real vector \(z_i\), the onsite quadratic form is at most
\(\eps(3|x_i|^2-k_i)\|z_i\|^2\), by Cauchy--Schwarz and \(\eps>0\). Since the largest eigenvalue of \(-gL\) is \(3g\), on every point of \eqref{eq:segment},
\begin{equation}
\nabla^2\V\preceq
\left[\frac1{10}\left(3\frac{121}{100}-\frac{999}{1000}\right)
+\frac3{60}\right]I_6
=\frac{3131}{10000}I_6.
\label{eq:curvature}
\end{equation}
Only an upper Hessian bound is needed; no positive lower bound on the disk product is asserted.

Taylor's integral formula and \(b=-\nabla\V(\Omega)\) now give
\begin{align}
\V(v)-\V(\Omega)
&=-\|b\|_2^2+
\int_0^1(1-t)b^T\nabla^2\V(\Omega+tb)b\,dt\notag\\
&\le-\left(1-\frac12\frac{3131}{10000}\right)\|b\|_2^2
=-\frac{16869}{20000}\|b\|_2^2.
\label{eq:predecrease}
\end{align}

\paragraph{Synchronization and the equality states.}
The prestage is nonzero by \eqref{eq:prebounds}, its phase width is \(D_v\le D_0<\pi/6\), and \(\lambda=1/300\) lies in the phase lemma's interval. The factor in \eqref{eq:phasebound} obeys
\[
0\le1-9\lambda\cos3D_v
\le1-\frac{3\sqrt2}{200}<1.
\]
For the fixed prestage moduli \(\rho_i=|v_i|\),
\begin{equation}
|\rho_i e^{i\psi_i}-\rho_j e^{i\psi_j}|^2
=(\rho_i-\rho_j)^2+
2\rho_i\rho_j[1-\cos(\psi_i-\psi_j)].
\label{eq:pairdistance}
\end{equation}
The function \(1-\cos t\) increases with \(|t|\) on this interval, so each pair distance is nonincreasing under the phase stage. Each onsite term is unchanged, and \(g>0\). Hence
\begin{equation}
\V(F\Omega)\le\V(v).
\label{eq:syncdecrease}
\end{equation}
Combining with \eqref{eq:predecrease} proves \eqref{eq:finitedescent}, including the domain's closed boundaries. This uses the native first-harmonic graph distances, not merely the separate harmonic-three phase scalar.

If \(\V(F\Omega)=\V(\Omega)\), the positive coefficient in \eqref{eq:finitedescent} forces \(b=0\). Write \(\Omega_i=r_i e^{i\theta_i}\) in a consistent interval and choose a channel \(m\) of minimal phase. The imaginary part of \(e^{-i\theta_m}b_m=0\) is
\begin{equation}
0=g\sum_{j\ne m}r_j\sin(\theta_j-\theta_m).
\label{eq:minimumphase}
\end{equation}
Every summand is nonnegative, and any strictly larger phase makes its summand positive because the difference is at most \(D_0\). Thus all phases agree. The remaining real equations are \eqref{eq:fixedbranch} in the amplitude cube, whose unique solution is the inherited \(r_*\). Conversely, \(e^{i\theta}r_*\) has \(v=\Omega\) and zero kicks. Therefore
\begin{equation}
\V(F\Omega)=\V(\Omega)
\quad\Longleftrightarrow\quad b=0
\quad\Longleftrightarrow\quad
\Omega\in\{e^{i\theta}r_*\},\qquad \Omega\in\D.
\label{eq:finiteequality}
\end{equation}
\end{proof}

The coefficient in \eqref{eq:finitedescent} is neither a percentage reduction of potential nor a bound using \(\|F\Omega-\Omega\|^2\). The result concerns the ideal mathematical map, not a certification of floating-point evaluation. It does not classify equality in every intermediate quantitative estimate.
